\documentclass[10pt,a4paper]{amsart}
\usepackage[margin=19mm]{geometry}
\usepackage{amsmath,amssymb,mathtools}
\usepackage[scr=rsfs,cal=euler]{mathalfa}
\usepackage{xfrac}
\usepackage[unicode,hidelinks,bookmarks=false]{hyperref}
\newcommand{\ie}{i.e.\ }
\newcommand{\cf}{cf.\ }
\newcommand{\resp}{respectively\ }
\newcommand{\Subsection}{\subsection}
\providecommand{\nobreakdash}{\nobreak\hbox{-}}
\setlength{\emergencystretch}{2em}
\multlinegap0pt
\usepackage{tcolorbox}
\usepackage{tikz}
\newtheorem{theorem}{Theorem}
\title{Discrete trace formulas and holomorphic functional calculus for the adjacency matrix of regular graphs}
\author{Yulin Gong, Wenbo Li, Shiping Liu}
\date{Source: arXiv:2406.17505v3 (CC BY 4.0)}
\begin{document}
\maketitle

\noindent\emph{Source and licence.} Discrete trace formulas and holomorphic functional calculus for the adjacency matrix of regular graphs. Version arXiv:2406.17505v3 (CC BY 4.0), distributed by arXiv under the Creative Commons Attribution 4.0 International licence, http://creativecommons.org/licenses/by/4.0/, as recorded in the arXiv licence metadata for this exact version and shown on the abstract page https://arxiv.org/abs/2406.17505v3. Full licence notice: \texttt{licenses/arxiv-2406.17505-CC-BY-4.0.txt}.\par\smallskip

\noindent\textit{Editorial scope.} Counted unit: the article's theorem master (source0.tex lines 568--581). The article's complete original proof of it is delivered with it (source0.tex lines 1254--1308, one \texttt{proof} environment of 55 lines, which the article places away from the statement). No mathematics is written, repaired or re-derived: the statement and the proof are the article's own lines, in the article's own notation, and no retained line is substituted or re-flowed.  Retained uncounted as the source-local context the proof needs: lines 307--309.  Nothing was omitted from any retained span, so the package carries no omission record.  No retained line contains a citation, so the wrapper carries no bibliography and no bibliography entry.  No retained line contains a figure environment, an image-inclusion command or drawing code, so no image is delivered and the package declares no asset dependency.  Editorial formatting only, in addition: an AMS \texttt{amsart} preamble that restates the article's own declarations and macros used by the retained lines, namely the environments \texttt{theorem} on separate counters, together with the standard packages those lines need.  The retained environments print in this document as Theorem 1 in order of appearance, whereas the article prints the same items with the numbering of its own full document.  The complete native source of the article as distributed in the arXiv e-print for this exact version is bundled unchanged as \texttt{sources/161346/source0.tex}.
    \begin{equation}\label{Xrqgeneratingfunction}
    \frac{1-q^{-1}t^2}{1-xt+t^2}=\sum_{r=0}^{\infty}X_{r,q}(x)t^{r}.
    \end{equation}

\begin{theorem}\label{master}
    Suppose that $h$ is holomorphic on $\Omega(\rho), \rho>0, \rho\neq 1$.
    Then for all $q\in \mathbb{N}\cup \{\infty\}$ and $z \in \Omega(\rho)$, we have
    \begin{equation}\label{masterlemma}
        h(z)=\sum_{r=0}^{\infty}a_{r,q}(h)X_{r,q}(z),
    \end{equation}
    where
    \begin{equation}\label{crq}
    \begin{aligned}
        a_{r,q}(h)=\frac{1}{2\pi i}\oint_{\partial B(0,1)} h(\xi+\xi^{-1})\xi^{r-1}\frac{1-\xi^2}{1-q^{-1}\xi^2}d\xi.
    \end{aligned}
    \end{equation}
    Moreover, the right hand side of (\ref{masterlemma}) converges absolutely and uniformly on any compact subset of $\Omega(\rho)$.
\end{theorem}

\begin{proof}[Proof of Theorem \ref{master}]
    Fix $z \in \Omega(\rho)$. Choose $0<\varepsilon<\rho-1$ such that $z \in \Omega(\rho-\varepsilon)$. By the Cauchy's integral formula, we have 
    $$h(z)=\frac{1}{2\pi i}\oint_{\partial \Omega(\rho-\varepsilon)}\frac{h(\omega)}{\omega-z}d\omega.$$
    By a change of variable $\omega=\xi+\xi^{-1}$, we have
    \begin{equation}
        \begin{aligned}
    h(z)&=\frac{1}{2\pi i}\oint_{\partial \Omega(\rho-\varepsilon)}\frac{h(\omega)}{\omega-z}d\omega\\
    &=\frac{1}{2\pi i}\oint_{\partial B(0, \frac{1}{\rho-\varepsilon})}\frac{h(\xi+\xi^{-1})}{\xi+\xi^{-1}-z}(\frac{1}{\xi^2}-1)d\xi\\
    &=\frac{1}{2\pi i}\oint_{\partial B(0, \frac{1}{\rho-\varepsilon})} \frac{1-\xi^2}{\xi(1-q^{-1}\xi^2)} h(\xi+\xi^{-1}) \frac{1-q^{-1}\xi^2}{1-z\xi+\xi^2}d\xi.
\end{aligned}
    \end{equation}
    Notice that by the choice of $\varepsilon$, 
    $$\frac{1-q^{-1}\xi^2}{1-z\xi+\xi^2}$$ has no poles on $\overline{ B(0, \frac{1}{\rho-\varepsilon})}$ as a function of $\xi$. By (\ref{Xrqgeneratingfunction}), 
    $$\frac{1-q^{-1}\xi^2}{1-z\xi+\xi^2}=\sum_{r=0}^{\infty}X_{r,q}(z)\xi^r,$$
    and the right hand side converges uniformly on $\overline{ B(0, \frac{1}{\rho-\varepsilon})}$. Thus 
    $$
    \begin{aligned}
    h(z)&=\frac{1}{2\pi i}\oint_{\partial B(0, \frac{1}{\rho-\varepsilon})} \frac{1-\xi^2}{\xi(1-q^{-1}\xi^2)} h(\xi+\xi^{-1}) \sum_{r=0}^{\infty}X_{r,q}(z) \xi^r d\xi\\
    &=\sum_{r=0}^{\infty} \left( \frac{1}{2\pi i}\oint_{\partial B(0, \frac{1}{\rho-\varepsilon})} h(\xi+\xi^{-1})\xi^{r-1}\frac{1-\xi^2}{1-q^{-1}\xi^2}d\xi \right) X_{r,q}(z).
\end{aligned}$$
Here we can interchange the order of integral and summation due to the uniform convergence mentioned above.
For $q>1$, the function 
$$h(\xi+\xi^{-1})\xi^{r-1}\frac{1-\xi^2}{1-q^{-1}\xi^2}$$ is holomorphic on $\{z:\rho^{-1}<|z|<\mathrm{min}(\rho, q^{1/2}) \}$. Thus 
$$
\begin{aligned}
    &\frac{1}{2\pi i}\oint_{\partial B(0, \frac{1}{\rho-\varepsilon})} h(\xi+\xi^{-1})\xi^{r-1}\frac{1-\xi^2}{1-q^{-1}\xi^2}d\xi\\
    &=\frac{1}{2\pi i}\oint_{\partial B(0, 1)} h(\xi+\xi^{-1})\xi^{r-1}\frac{1-\xi^2}{1-q^{-1}\xi^2}d\xi.
\end{aligned}
$$
For $q=1$, the coefficient of $X_{r,1}$ equals
$$\frac{1}{2\pi i}\oint_{\partial B(0, \frac{1}{\rho-\varepsilon})} h(\xi+\xi^{-1})\xi^{r-1}d\xi.$$
Similarly, we have 
$$
\begin{aligned}
    \frac{1}{2\pi i}\oint_{\partial B(0, \frac{1}{\rho-\varepsilon})} h(\xi+\xi^{-1})\xi^{r-1}d\xi=\frac{1}{2\pi i}\oint_{\partial B(0, 1)} h(\xi+\xi^{-1})\xi^{r-1}d\xi.
\end{aligned}
$$

Next, we prove the absolute and uniform convergence. Fix any compact subset $K\subset \Omega(\rho)$. There exits $0<\varepsilon_K<\rho-1$ such that $K\subset \Omega(\rho-\varepsilon_K)$. For any $z\in \Omega(\rho-\varepsilon_K)$, there exists $t\in \mathbb{C}$, such that $1\leq |t|\leq \rho-\varepsilon_K$ amd $z=t+t^{-1}$.
Since $$
    X_r(z)=\frac{t^{r+1}-t^{-r-1}}{t-t^{-1}}=\sum_{i=0}^{r}t^{r-2i},
$$
we have 
$$X_{r,q}(z)=X_r(z)-q^{-1}X_{r-2}(z)=t^{r}+t^{-r}+(1-q^{-1})\sum_{i=0}^{r}t^{r-2i}.$$
Hence
$$|X_{r,q}(z)|\leq |t|^{r}+|t|^{-r}+(1-q^{-1})\sum_{i=0}^{r}|t|^{r-2i}\leq (r+1)|t|^{r}\leq (r+1) (\rho-\varepsilon_K)^r.$$
For $q>1$, as shown above,
$$a_{r,q}(h)=\frac{1}{2\pi i}\oint_{\partial B(0, \tau)} h(\xi+\xi^{-1})\xi^{r-1}\frac{1-\xi^2}{1-q^{-1}\xi^2}d\xi,$$
for all $ \rho^{-1}<\tau\leq 1$. Thus 
$$|a_{r,q}|\leq M(\tau) \tau^r$$
where $$M(\tau)=\sup_{|\xi|=\tau}\ h(\xi+\xi^{-1})\frac{1-\xi^2}{1-q^{-1}\xi^2}.$$
Choose $\rho^{-1} <\tau<(\rho-\varepsilon_K)^{-1}$, there holds for any $z\in K$
$$\left|\sum_{r=0}^{\infty} a_{r,q}(h) X_{r,q}(z)\right|\leq \sum_{r=0}^{\infty} M(\tau) (r+1) \left(\frac{\tau}{\rho-\varepsilon_K}\right) ^r<+\infty.$$
The same proof can be applied to the case $q=1$. Thus the convergence is absolute and uniform on the compact subset $K$. 
\end{proof}

\end{document}
